\section{Threats to Validity and Limitations}
\label{sec:limitations}

\subsection{Construct and data validity}

NSRDB GHI is a satellite-informed physical-model estimate
\cite{sengupta2018}, not an on-site pyranometer record. Factory coordinates
select product grid cells, so neither the target nor its error quantifies the
irradiance at a particular roof surface. The purposive set of ten
factory-associated locations is also not a sample of factory electricity
demand, photovoltaic output, roof area, or operational constraints. The
industrial framing therefore identifies geographically relevant forecasting
sites; it does not establish plant-level energy savings.

The matched input deliberately contains only causal GHI and site identity.
This restriction isolates architectural effects, but it omits potentially
useful exogenous forecasts such as cloud motion, numerical weather prediction,
aerosols, and temperature. A site embedding can also encode persistent
location-specific behaviour and does not demonstrate transfer to an unseen
factory. Fixed local offsets produce regular hourly grids but do not reproduce
daylight-saving clock changes.

\subsection{Temporal and statistical validity}

All test evidence is retrospective and restricted to 2024. It does not measure
prospective robustness under later satellite-product revisions, climate
nonstationarity, or new sites. The monthly primary task has only 48 pre-test
refit origins per site, and the exploratory six-month task has seven test
origins per site. Its long output therefore cannot support the same empirical
resolution as the hourly or daily tasks.

Targets from adjacent multi-step origins overlap, creating strongly dependent
errors. Site-first aggregation avoids allowing that overlap to change spatial
weights, but it does not make the origins independent. The paired summaries
accordingly use ten locations as descriptive units and do not attach an
i.i.d. significance test to thousands of overlapping target points.

\subsection{Model-comparison validity}

TimesNet and DilatedRNN share inputs, origins, seeds, optimization limits, and
post-processing in the daily and monthly tasks, but their parameterizations and
inductive biases are not identical. The frozen configuration is a controlled
comparison rather than an exhaustive architecture-specific search. Moreover,
the hourly extension adds DilatedRNN with the single official seed 42 to
read-only comparator forecasts, whereas daily and monthly estimates use five
seeds. Cross-resolution claims therefore concern protocol behaviour and are
not interpreted as equal-precision estimates of random-seed variability.

The zero floor imposes a defensible physical bound on all resolutions. The
additional hourly nighttime mask uses deterministic solar geometry, but no
common upper physical envelope is imposed. Consequently, implausibly high
daytime extrapolations remain visible in the errors rather than being hidden by
post-processing.

\subsection{Contextual-data validity}

The K\"oppen--Geiger labels are categorical cells from the present nominal
1-km map \cite{beck2018}; they do not resolve microclimates within a cell or
conditions on a specific roof. Class boundaries are uncertain, as made explicit
by the 50\% confidence of the BMW San Luis Potos\'i assignment. Climate labels
are descriptive strata, not model inputs and not explanations of individual
forecast errors.

NASA POWER is likewise a gridded, model-based external product rather than
local rain-gauge or ceilometer evidence \cite{nasapowerdaily}. The 1~mm/day
precipitation and 80\% cloud thresholds provide reproducible post-hoc tags, but
thresholding cannot establish that weather caused a particular error.
Weather-case observations therefore supplement the full aggregate evaluation;
they do not replace it or determine which model is reported.
